\section{A normalized matrix product state in the final span}
\label{sec:lvv_final_mps}

The finite-span estimate and the full-left viable witness give the final extraction
in \cite{Landau2015Polynomial}. The argument uses the actual finite family and
does not appeal to a separate ground-state approximation theorem.

\begin{lemma}[Bond bound for a finite left span]
    \label{thm:lvv_left_span_bond}
    \lean{TNLean.MPS.AreaLaw.isOpenMPS_leftChain_of_mem_span}
    \leanok
    \uses{thm:left_mps_finite_span}
    Let $S$ be a finite family of vectors on the first $i$ sites. Suppose each
    member, read as a vector on an $i$-site chain, has MPS bond dimension at most
    $\max\{1,b\}$. Then every $v\in\operatorname{span}S$, in the same coordinates,
    is an MPS with bond dimension at most $|S|\max\{1,b\}$.
\end{lemma}
\begin{proof}
    \leanok
    Write $v=\sum_{s\in S}c_ss$. At each internal cut $k$, the coefficient
    matrix satisfies $M_k(v)=\sum_{s\in S}c_sM_k(s)$. Scalar multiplication
    does not increase matrix rank and rank is subadditive under finite sums.
    Therefore
    \begin{align}
        \operatorname{rank}M_k(v)
        \le\sum_{s\in S}\operatorname{rank}M_k(s)
        \le|S|\max\{1,b\}.
    \end{align}
    The characterization of open-boundary MPS by their internal cut ranks gives
    the assertion. Reading the left block as a chain is linear; if two members
    become equal under this reading, the image family has smaller cardinality.
\end{proof}

\begin{lemma}[Normalized MPS witness from a viable final family]
    \label{thm:lvv_final_span_mps}
    \lean{TNLean.MPS.AreaLaw.IsViable.exists_openMPS_of_forall_isLeftOpenMPS}
    \leanok
    \uses{thm:lvv_left_span_bond}
    Let $S$ be a finite $(L,\delta)$-viable family for a chain vector $\Gamma$.
    Suppose its members have individual effective MPS bond bound $\max\{1,b\}$.
    There exists a full-chain vector $\psi$ such that
    \begin{align}
        \|\psi\|=1,\qquad
        |\langle\psi,\Gamma\rangle|\ge1-\delta,
    \end{align}
    and $\psi$ is an open-boundary MPS of bond dimension at most
    $|S|\max\{1,b\}$.
\end{lemma}
\begin{proof}
    \leanok
    At the final cut the right block has a single configuration. Thus a unit
    viable witness, read on the full left block, lies in $\operatorname{span}S$.
    This reading is an isometry, so the witness retains its norm and overlap.
    Apply the finite-span bond estimate.
\end{proof}

\begin{lemma}[Final extraction with the recorded size bound]
    \label{thm:lvv_final_viable_mps}
    \lean{TNLean.MPS.AreaLaw.IsViableMPS.exists_openMPS}
    \leanok
    \uses{thm:lvv_final_span_mps}
    An $(L,s,b,\delta)$-viable family gives a normalized full-chain MPS
    $\psi$ of bond dimension at most $s\max\{1,b\}$ with
    $|\langle\psi,\Gamma\rangle|\ge1-\delta$.
\end{lemma}
\begin{proof}
    \leanok
    Its actual cardinality is at most $s$. Enlarge the finite-span bond bound
    by this inequality.
\end{proof}

\begin{theorem}[Polynomial final bond at prescribed accuracy]
    \label{thm:lvv_sampled_ground_state}
    \lean{TNLean.MPS.AreaLaw.exists_sampled_groundState_openMPS_polynomial}
    \leanok
    For fixed $q\ge2$, $\Delta>0$, $\eta>0$, an integer $a\ge1$, and $B$, there are natural
    constants $C\ge1,p$ such that every positive nearest-neighbour
    contraction chain of length $L\ge1$, with a unique gapped ground vector
    $\Gamma$ and $\varepsilon_0/\Delta\le B$, admits a normalized matrix
    product state $\psi$ of bond at most $CL^p$ with
    $|\langle\psi,\Gamma\rangle|\ge1-\eta/L^a$.
    The construction is the successive sampled construction.
\end{theorem}
\begin{proof}
    \leanok
    \uses{thm:lvv_sampled_successive_families, thm:lvv_final_span_mps}
    Apply Lemma~\ref{thm:lvv_final_span_mps} to $S_L$.
    Multiply the size and bond polynomials. The overlap error is at most
    $\alpha/L^a\le\eta/L^a$.
\end{proof}
